\documentclass[../main.tex]{subfiles}
\begin{document}
\chapter{Principle of Least Action}
\section{Configuration Space}
Consider a system of $N$ particles in 3D with positions $\vec{x}_1, \ldots, \vec{x}_n$.
We can consider the positions of the particles at a single point as a $3N$ dimensional point in \textit{configuration space}:
\[
  (\vec{x}_1, \ldots, \vec{x}_N)
\]
It is sometimes also useful to use non-Cartesian coordinates for the configuration space:
\[
  \vec{q} = (\vec{q}_1, \ldots, \vec{q}_n)
\]
\section{The Lagrangian}
The \textit{principle of least action} gives us a simple way to determine the equations of motion from the kinetic energy $T$ and potential energy $V$ in terms of $\vec{q}(t)$ and $\dot{\vec{q}}(t)$.
\begin{definition}[Lagrangian]
  We define the \textit{Lagrangian} as:
  \[
    L(\vec{q}, \dot{\vec{q}}, t) = T - V
  \]
\end{definition}
\begin{definition}[Action]
  The \textit{action} of a path $\vec{q}(t)$ in configuration space between $\vec{q}_A$ and time $t_A$, and $\vec{q}_B$ at time $t_B$ is given by the functional:
  \[
    I[\vec{q}] = \int_{t_A}^{t_B} L(\vec{q}(t), \dot{\vec{q}}(t), t) \d{t}
  \]
\end{definition}
\begin{law}[Principle of Least Action]
  The \textit{principle of least action} tells us that the actual path from $\vec{q}_A$ to $\vec{q}_B$ of the system extremises the action.
\end{law}
Since $\vec{q}_A$ and $\vec{q}_B$ are fixed, the boundary term in $\delta I$ vanishes and so we can utilise the Euler-Lagrange equations.
As $\vec{q}$ has $3N$ components, we will obtain $3N$ Euler-Lagrange equations:
\[
  \pderiv{L}{q_i} - \deriv{}{t} \pderiv{L}{\dot{q}_i} = 0 \text{ for } i \in \{1, \ldots, 3N\}
\]
\begin{example}
  If we take $N = 1$ and $\vec{q} = \vec{x}$ (i.e. we work in Cartesian coordinates), we know that $T = \frac{1}{2} m \dot{\vec{x}} \cdot \dot{\vec{x}}$ for fixed $m$ and assume that we can write $V = V(\vec{x}, t)$. \label{n1LG}
  The Lagrangian is then:
  \[
    L = \frac{1}{2} m \dot{\vec{x}} \cdot \dot{\vec{x}} - V(\vec{x}, t)
  \]
  and so the Euler-Lagrange equations tells us that:
  \[
    -\pderiv{V}{x_i} - \deriv{}{t}(m \dot{x}_i) = 0 \ \forall i \in \{1, 2, 3\} \iff m \ddot{\vec{x}} = -\nabla V
  \]
  This is Newton's 2nd Law for a conservative force $\vec{F} = - \nabla V$.
\end{example}
\begin{proposition}
  If $L$ has no explicit $t$ dependence, then for any solution of the Euler-Lagrange equations there exists a first integral: \label{LnoT}
  \[
    L - \sum_{i = 1}^{3N}  \dot{q}_i \pderiv{L}{\dot{q}_i} = - E
  \]
  for some constant $E$.
\end{proposition}
\begin{proof}
  We first consider the full derivative $\deriv{L}{t}$ which is not necessarily 0 because of the implicit $t$ dependence:
  \begin{align*}
    \deriv{L}{t} &= \deriv{}{t} L(\vec{q}(t), \dot{\vec{q}}(t), t) \\
                 &= \cancelto{0}{\pderiv{L}{t}} + \sum_{i = 1}^{3N} \left(\pderiv{L}{q_i} \dot{q}_i + \pderiv{L}{\dot{q_i}} \ddot{q}_i\right) \text{ by multivariate chain rule}\\
                 &= \sum_{i = 1}^{3N} \dot{q}_i \cancelto{0}{\left(\pderiv{L}{q_i} - \deriv{}{t} \pderiv{L}{\dot{q_i}}\right)} + \sum_{i = 1}^{3N} \left(\ddot{q}_i \pderiv{L}{\dot{q_i}} + \dot{q}_i \deriv{}{t} \pderiv{L}{\dot{q_i}}\right) \\
                 &= \deriv{}{t}\left(\sum_{i = 1}^{3N} \dot{q}_i \pderiv{L}{\dot{q_i}}\right)
  \end{align*}
  Rearranging we have:
  \[
    \deriv{}{t}\left(L - \sum_{i = 1}^{3N} \dot{q}_i \pderiv{L}{\dot{q}_i}\right) = 0 \implies L - \sum_{i = 1}^{3N} \dot{q}_i \pderiv{L}{\dot{q}_i} = -E
  \]
  for some constant $E$.
\end{proof}
\begin{example}
  Continuing from \cref{n1LG}, if we further assume that $V$ has no $t$ dependence, i.e. $V = V(\vec{x})$, then the constant $E$ is:
  \begin{align*}
    -E &= \frac{1}{2} m \dot{\vec{x}} \cdot \dot{\vec{x}} - V - \sum_{i = 1}^{3} \dot{x}_i m \dot{x}_i \\
       &= -\frac{1}{2} m \dot{\vec{x}} \cdot \dot{\vec{x}} - V(\vec{x}) \\
       &= - (T + V)
  \end{align*}
  So $E = T + V$ is the \textbf{conserved} total energy.
\end{example}
\begin{remark}
  Because we can recover the equations of motion from the Lagrangian, we think of it as \textit{encoding} the laws of physics.
  So we can think of $L$ having no $t$ dependence as being the same as the statement that the laws of physics are time translation invariant.
\end{remark}
\begin{example}[Polar Coordinates]
  Consider a particle in 2D dimensions with position $\vec{q} = (r, \theta)$ in plane polar coordinates.
  In polar coordinates, the kinetic energy $T$ is given by:
  \[
    T = \frac{1}{2}m(\dot{r}^2 + r^2 \dot{\theta}^2)
  \]
  If we assume that the force is radial, i.e. $V = V(r)$, then the Lagrangian is:
  \[
    L = \frac{1}{2}m(\dot{r}^2 + r^2 \dot{\theta}^2) - V(r)
  \]
  Using the Euler-Lagrange equation for $\theta(t)$:
  \[
    \pderiv{L}{\theta} - \deriv{}{t} \pderiv{L}{\dot{\theta}} = 0
  \]
  We see that $\pderiv{L}{\theta} = 0$ so we have a first integral:
  \[
    \pderiv{L}{\dot{\theta}} = mr^2 \dot{\theta} = \text{constant}
  \]
  which is the conservation of angular momentum.
  Since $m$ is constant, we call the conserved angular momentum per unit mass $h = r^2 \dot{\theta}$.

  We could try and use the Euler-Lagrange equation for $r(t)$, however there is explicit $r$ dependence and so we would end up with a second order ODE.
  However, since $\pderiv{L}{t} = 0$, we know from \cref{LnoT} that there is a conserved energy:
  \[
    E = \dot{r} \pderiv{L}{\dot{r}} + \dot{\theta} \pderiv{L}{\dot{\theta}} - L = \frac{1}{2}m (\dot{r}^2 + r^2 \dot{\theta}^2) + V = T + V
  \]
  Solving to eliminate $\dot{\theta}$, we obtain:
  \[
    \frac{1}{2} m\dot{r}^2 + V_{\text{eff}}(r) = E
  \]
  where:
  \[
    V_{\text{eff}}(r) = \frac{mh^2}{2r^2} + V(r)
  \]
  We can then go on to analyse this using methods as discussed in IA Dynamics \& Relativity.
\end{example}
\section{The Hamiltonian}
\subsection{Definition}
\begin{definition}
  We assume that the Lagrangian $L(\vec{q}, \dot{\vec{q}}, t)$ is a convex function of $\dot{\vec{q}}$.
  We then define the \textit{Hamiltonian} as the Legendre transform of $L$ with respect to $\dot{\vec{q}}$:
  \begin{align*}
    H(\vec{q}, \vec{p}, t) &= \sup_{\dot{\vec{q}}} [\vec{p} \cdot \dot{\vec{q}} - L(\vec{q}, \dot{\vec{q}}, t)] \\
                           &= \at{[\vec{p} \cdot \dot{\vec{q}} - L(\vec{q}, \dot{\vec{q}}, t)]}{\dot{\vec{q}} = \dot{\vec{q}}(\vec{q}, \vec{p}, t)}
  \end{align*}
  where $\dot{\vec{q}}(\vec{q}, \vec{p}, t)$ is defined to be the local extremum, which by convexity must exist and also be the global extremum.
  We also determine $\vec{p}$ from:
  \[
    p_i = \pderiv{L}{\dot{q}_i}
  \]
\end{definition}
\subsection{Examples}
\begin{example}
  Consider a single particle (i.e. $N = 1$) with $\vec{q} = \vec{x}$.
  \[
    L = \frac{1}{2} m \dot{\vec{x}} \cdot \dot{\vec{x}} - V(\vec{x}, t)
  \]
  We then have:
  \[
    p_i = \pderiv{L}{\dot{x}_i} = m \dot{x}_i \implies \vec{p} = m \dot{\vec{x}}
  \]
  which is the momentum, and:
  \[
    H = \vec{p} \cdot \frac{\vec{p}}{m} - \left(\frac{\vec{p} \cdot \vec{p}}{2m} - V\right) = \frac{\vec{p} \cdot \vec{p}}{2m} + V
  \]
  which is the total energy.
\end{example}
\begin{example}
  Consider a single particle with a central force where $\vec{q} = (r, \theta)$, we know that:
  \[
    L = \frac{1}{2} (\dot{r}^2 + r^2 \dot{\theta}^2) - V(r)
  \]
  so:
  \[
    p_r = \pderiv{L}{\dot{r}} = m\dot{r} \text{ and } p_{\theta} = \pderiv{L}{\dot{\theta}} = mr^2 \dot{\theta}
  \]
  and:
  \[
    H = p_r \dot{r} + p_{\theta} \dot{\theta} - L = \frac{p^2_r}{2m} + \frac{p^{2}_{\theta}}{2m} + V(r)
  \]
\end{example}
In general, we call $p_i$ the \textit{conjugate momentum} to $q_i$.

After we have eliminated $\dot{\vec{q}}$ in favor of $\vec{q}, \vec{p}, t$ by using $p_i = \pderiv{L}{\dot{q}_i}$, we have:
\[
  H = \sum_{i} \dot{q}_i \pderiv{L}{\dot{q}_i} - L
\]
so $H$ is the total energy and so will be conserved if $\pderiv{L}{t} = 0$.
\subsection{Hamiltonian Equations of Motion}
We can writ the equations of motion in terms of $H$:
\[
  \at{\pderiv{H}{p_i}}{\vec{q}} = \dot{q}_i + \sum_{j} p_j \pderiv{\dot{q}_j}{p_i} - \sum_{j} \pderiv{L}{\dot{q}_j} \pderiv{\dot{q}_j}{p_i} = \dot{q}_i
\]
and:
\begin{align*}
  \at{\pderiv{H}{q_i}}{\vec{p}} &= \sum_{j} p_j \pderiv{\dot{q}_j}{q_i} - \pderiv{L}{q_i} - \pderiv{L}{q_i} - \sum_{j} \pderiv{L}{\dot{q}_j} \pderiv{\dot{q}_j}{q_i} \\
                                &= - \pderiv{L}{q_i} \\
                                &= - \deriv{}{t} \pderiv{L}{\dot{q}_i} \text{ using the Euler-Lagrange equations} \\
                                &= - \dot{p}_i
\end{align*}
Thus the Euler Lagrange equations imply \textit{Hamilton's Equations} :
\[
  \dot{q}_i = \pderiv{H}{p_i} \text{ and } \dot{p}_i = - \pderiv{H}{q_i}
\]
The Euler-Lagrange equations are second order ODEs in a $3N$ dimensional \textit{configuration space} $\vec{q}$, whereas Hamilton's equations are first order ODEs in a $6N$ dimensional \textit{phase space} with coordinates $(\vec{q}, \vec{p})$.

We can in fact also obtain these equations as the Euler-Lagrange equations for the functional:
\[
  I[\vec{q}, \vec{p}] = \int_{t_A}^{t_B} (\dot{\vec{q}} \cdot \vec{p} - H(\vec{q}, \vec{p}, t)) \d{t}
\]
for fixed $\vec{q}(t_A)$ and $\vec{q}(t_B)$ so that $\delta \vec{q}(t_A) = \delta \vec{q}(t_B) = 0$.

The time derivative of the Hamiltonian is:
\begin{align*}
  \deriv{H}{T} &= \deriv{}{t}H(\vec{q}(t), \vec{p}(t), t) \\
               &= \sum_{i} \left(\pderiv{H}{q_i}\dot{q}_i + \pderiv{H}{p_i}\dot{p}_i\right) + \pderiv{H}{t} \\
               &= \sum_{i} (-\dot{p}_i \dot{q}_i + \dot{q}_i \dot{p}_i) + \pderiv{H}{t} \\
               &= \pderiv{H}{t}
\end{align*}
Thus if $H$ has \textbf{no explicit time dependence}, then the Hamiltonian is conserved.
\section{Symmetry and Noether's Theorem}
We want to explore the idea that continuous symmetries of the action are in correspondence the 1st integrals of the Euler-Lagrange equations.
\subsection{Symmetries of the Action}
Consider a smooth map $\vec{Q}: \R^{6N + 1} \to \R^{3N}$.
Given a curve $\vec{q}(t)$, we can define the curve $\vec{q}^{*}$:
\[
  \vec{q}^{*}(t) = \vec{Q}(\vec{q}(t), \dot{\vec{q}}(t), t)
\]
\begin{definition}[Symmetry of the Action]
  We say that $\vec{Q}$  is a \textit{symmetry of the action} if for all curves $\vec{q}(t)$ and all $t_A, t_B$ we have:
  \[
    \int_{t_A}^{t_B} L(\vec{q}^{*}(t), \dot{\vec{q}}^{*}(t), t) \d{t} = \int_{t_A}^{t_B} L(\vec{q}(t), \dot{\vec{q}}(t), t) \d{t}
  \]
\end{definition}
Now consider a \textit{1-parameter family} of such maps, that is we have a collection of curves characterised by the parameter $s$:
\[
  \vec{Q}(\vec{q}, \dot{\vec{q}}, t, s)
\]
with $\vec{Q}(\vec{q}, \dot{\vec{q}}, t, 0) = \vec{q}$.
For small $s$:
\[
  \vec{q}^{*} = \vec{q} + s \at{\pderiv{\vec{Q}}{s}}{s =0} + O(s^2)
\]
\begin{theorem}[Noether's Theorem -- V1]
  If there exists a 1-parameter family of symmetries of the action then any solution of the Euler-Lagrange equations has a conserved quantity (i.e. a first integral):
  \[
    \sum_{i = 1}^{3N} \pderiv{L}{\dot{q}_i} \at{\pderiv{Q_i}{s}}{s = 0} = \text{constant}
  \]
\end{theorem}
\begin{proof}
  For small $s$, we obtain:
  \begin{align*}
    &\int_{t_a}^{t_B} \left[L\left(\vec{q}(t) + s \at{\pderiv{\vec{Q}}{s}}{s = 0} + \cdots,  \dot{\vec{q}}(t) + s \deriv{}{t}\left(\at{\pderiv{\vec{Q}}{s}}{s = 0}\right) + \cdots\right) -L(\vec{q}(t), \dot{\vec{q}}(t), t)\right] \d{t} \\
    &= \int_{t_A}^{t_B} \left[s \at{\pderiv{Q_i}{s}}{s=0} \pderiv{L}{q_i} + s \deriv{}{t}\left(\at{\pderiv{Q_i}{s}}{s = 0}\right) \pderiv{L}{\dot{q}_i} + \cdots\right] \d{t} \\
    &=\int_{t_A}^{t_B} \left[s \at{\pderiv{Q_i}{s}}{s = 0} \underbrace{\left(\pderiv{L}{q_i} - \deriv{}{t}\pderiv{L}{\dot{q}_i}\right)}_{=0} + s \deriv{}{t}\left(\pderiv{L}{\dot{q}_i} \at{\pderiv{Q_i}{s}}{s = 0}\right) + \cdots\right] \d{t} \\
    &= \eval{s \pderiv{L}{q_i} \at{\pderiv{Q_i}{s}}{s = 0}}{t_A}{t_B} + O(s^2)
  \end{align*}
  But this whole integral is equal to 0 and so $\pderiv{L}{\dot{q}_i} \at{\pderiv{Q_i}{s}}{s = 0}$ takes some value at $t_A, t_B$, but $t_A$ and $t_B$ were arbitrary and so we are done.
\end{proof}
\begin{example}[Space Translational Symmetry in One Coordinate]
  If $L$ does not depend on, for example $q_1$, then consider:
  \[
    \vec{Q} = (q_1 + s, q_2, q_3, \ldots, q_{3N})
  \]
  Thus:
  \begin{align*}
    q^{*}_1 &= q_1 + s \text{ and } q^{*}_i = q_i\ \forall i \in {2, \ldots, 3N} \\
    \dot{q}^{*}_i &= q_i\ \forall i \in {1, 2, \ldots, 3N}
  \end{align*}
  This is a symmetry of the action since $L$ is independent of $q_1$, thus we have a conserved quantity:
  \[
    \at{\pderiv{Q_1}{s}}{s = 0} \pderiv{L}{\dot{q}_1} = \pderiv{L}{\dot{q_1}} = p_1
  \]
\end{example}
\begin{example}[Translational Symmetry in Euclian Space]
  Consider Cartesian Coordinates  $\vec{q} = (\vec{x}_1, \ldots, \vec{x}_N)$ and assume $L$ depends on $\vec{x}_i$ only by $\vec{x}_i - \vec{x}_j$.
  Let $\vec{Q} = (\vec{x}_1 + s\vec{a}, \vec{x}_2 + s\vec{a}, \ldots, \vec{x}_N + s\vec{a})$.
  We see that $\dot{\vec{x}}_i - \dot{\vec{x}}_j = \vec{x}_i - \vec{x}_j$ and so $L(\vec{q}^{*}, \dot{\vec{q}}^{*}, t) = L(\vec{q}, \dot{\vec{q}}, t)$ and so we can apply Noether's theorem to conclude that the quantity:
  \[
    \sum_{i} \pderiv{L}{\dot{q_i}} \at{\pderiv{Q_i}{s}}{s = 0} = \vec{a} \cdot \sum_{i} \pderiv{L}{\dot{\vec{x}}_i} = \vec{a} \cdot \sum_{i} \vec{p}_i = \vec{a} \cdot \vec{P}
  \]
  is conserved where $\vec{P}$ is the total momentum of the system.
  Since $\vec{a}$ was arbitrary, we conclude that $\vec{P}$ is conserved.
  Thus the invariance of physical laws under spatial transformations implies the conservation of momentum.
\end{example}
\begin{example}[Rotational Symmetries of Cartesian Coordinates]
  Consider $\vec{q} = (\vec{x}_1, \ldots, \vec{x}_N)$.
  A rotation $R$ maps:
  \[
    \vec{q} \mapsto \vec{q}_{R} = (R\vec{x}_1, \ldots, R\vec{x}_N)
  \]
  If we assume that $L$ is invariant under notations then:
  \[
    L(\vec{q}_R, \dot{\vec{q}}_R, t) = L(\vec{q}, \dot{\vec{q}}, t) \text{ for all rotations $R$}
  \]
  Taking $Q = \vec{q}_{R(s)}$ where $R(s)$ is a rotation through angle $s$ about some axis $\uvec{n}$.
  For infinitesimal $s$:
  \[
    R(s)\vec{x} = \vec{x} + s \vec{n} \times \vec{x}
  \]
  and so:
  \[
    \at{\pderiv{Q}{s}}{s = 0} = (\uvec{n} \times \vec{x}_1, \ldots, \uvec{n} \times \vec{x}_N)
  \]
  So we have a conserved quantity:
  \begin{align*}
    \sum_{i, a} \pderiv{L}{(\dot{\vec{x}}_i)_a} (\uvec{n} \times \vec{x}_i)_{a} &= \sum_{i} \vec{p}_i \cdot (\uvec{n} \times \vec{x}_i) \text{ where $a \in \{1, 2, 3\}$} \\
                                                                                &= \vec{n} \cdot \sum_{i} \vec{x}_i \times \vec{p}_i \\
                                                                                &= \uvec{n} \cdot \vec{L}
  \end{align*}
  where $\vec{L}$ is the total angular momentum of the system.
  Since the axis $\uvec{n}$ was arbitrary, we conclude that $\vec{L}$ is conserved.
  Thus the invariance of physical laws under rotation implies the conservation of angular momentum.
\end{example}
We can generalise the definition of a 1-parameter family of symmetries to allow a boundary term:
\[
  \int_{t_A}^{t_B} L(\vec{q}^{*}, \dot{\vec{q}}^{*}, t) \d{t} = \int_{t_A}^{t_B} L(\vec{q}, \dot{\vec{q}}, t) \d{t} + \eval{sK(\vec{q}, \dot{\vec{q}}, t)}{t_A}{t_B}
\]
for small $s$ and some $K$.
It can be checked that even with this generalised definition, Noether's theorem still works, but the conserved quantity is now:
\[
  \sum_{i} \pderiv{L}{\dot{q}_i} \at{\pderiv{Q_i}{s}}{s = 0} - K
\]
\begin{example}
  \[
    L = \sum_{i = 1}^{N} \frac{1}{2} m_i \dot{\vec{x}}_i \cdot \dot{\vec{x}}_i - V(\vec{x}_1, \ldots, \vec{x}_N)
  \]
  where $V$ depends only on the separations $\vec{x}_ - \vec{x}_j$.
  From IA Dynamics and Relativity, we know a Galilean Boost is of the form $\vec{x}^{*} = \vec{x} - \vec{v}t$, we define:
  \[
    Q = (\vec{x} - \uvec{n} st, \vec{x}_2 - \uvec{n} st, \ldots)
  \]
  for $\vec{v} = \uvec{n} s$ and $\vec{x}^{*}_i = \vec{x}_i - \uvec{n}st$, $\dot{\vec{x}}^{*}_i = \dot{\vec{x}}_i - \uvec{n} s$.
  We then have:
  \begin{align*}
    L(\vec{x}^{*}, \dot{\vec{x}}^{*}, t) &= \sum_{i = 1}^{N} \frac{1}{2} m_i (\dot{\vec{x}}_i - \uvec{n} s) \cdot (\dot{\vec{x}}_i - \uvec{n} s)  - V(\vec{x}^{*}_1, \ldots, \vec{x}^{*}_N) \\
                                         &= \sum_{i = 1}^{N} \frac{1}{2} m_i \dot{\vec{x}}_i \cdot \dot{\vec{x}}_i - V(\vec{x}_1, \ldots, \vec{x}_N) - s \uvec{n} \sum_{i} m_i \dot{\vec{x}}_i + O(s^2) \\
                                         &= L(\vec{x}, \dot{\vec{x}}, t) - \deriv{}{t}(s \uvec{n} \cdot \sum_{i} m_i \vec{x}_i)
  \end{align*}
  since $\vec{x}^{*}_i - \vec{x}^{*}_j = \vec{x}_i - \vec{x}_j$.

  This is a 1-parameter symmetry but with boundary term:
  \[
    K = - \vec{n} \cdot \sum_{i} m_i \vec{x}_i = - \vec{n} \cdot (M \vec{R})
  \]
  where $M$ is the total mass and $\vec{R}$ is the centre of mass.
  Thus we have a conserved quantity:
  \begin{align*}
    \sum_{i} \pderiv{L}{\dot{q}_i} \at{\pderiv{Q_i}{s}}{s = 0} - K &= \sum_{i} \vec{p}_i \cdot (- \uvec{n} t) + M \uvec{n} \cdot \vec{R} \\
                                                                   &= \uvec{n} \cdot (M\vec{R} - \vec{P} t)
  \end{align*}
  Since $\uvec{n}$ was arbitrary, the quantity $M\vec{R} - \vec{P}t $ is conserved.
\end{example}
A different type of symmetry of the action involves transforming the time coordinate $t^{*} = T(t)$ for an invertible function $T$ and set $q^{*}(t^{*}) = \vec{q}(t) = \vec{q}(T^{-1}(t^{*}))$.
This is a symmetry of the action if:
\begin{equation}
  \int_{t^{*}_A}^{t^{*}_B} L(\vec{q}^{*}(t^{*}), \deriv{\vec{q}^{*}}{{t^{*}}}(t^{*}), t^{*}) \d{t^{*}} = \int_{t_{A}}^{t_{B}} L(\vec{q}(t), \dot{\vec{q}}(t), t) \d{t} \label{timeSym}
\end{equation}
for all $\vec{q}(t)$ and $t_A, t_B$.

Consider a 1-parameter family $T(t, s)$ where $T(t, 0) = t$:
\[
  t^{*} = t + s \at{\pderiv{T}{s}}{s = 0} + O(s^2)
\]
\begin{theorem}[Noether's Theorem V2]
  If there exists a 1-parameter family of symmetries of the action of the type defined about then any solution of the Euler-Lagrange equation has a conserved quantity:
  \[
    (L - \sum_{i} \dot{q}_i \pderiv{L}{\dot{q_i}})\at{\pderiv{T}{s}}{s = 0} = \text{constant}
  \]
\end{theorem}
\begin{proof}
  First note that:
  \[
    \deriv{{t^{*}}}{t} = 1 + s \deriv{}{t}\at{\pderiv{T}{s}}{s = 0} + O(s^2)
  \]
  Let $I^{*}$ and $I$ be the left hand sides and right hand sides of \cref{timeSym}.
  We change variables in $I^{*}$ from $t^{*}$ to $t$:
  \[
    I^{*} = \int_{t_A}^{t_B} L\left(\vec{q}(t), \deriv{t}{{t^{*}}} \deriv{}{t}\vec{q}(t), T(t)\right) \deriv{t^{*}}{t} \d{t}
  \]
  We can find:
  \[
    \deriv{t}{{t^{*}}} = \left(\deriv{t^{*}}{t}\right)^{-1} = 1 - s \deriv{}{t} \at{\pderiv{T}{s}}{s = 0} + O(s^2)
  \]
  and so:
  \begin{align*}
    I^{*} &= \int_{t_A}^{t_B} L\left(\vec{q}(t), \dot{\vec{q}}(t) - s \deriv{}{t}\at{\pderiv{T}{s}}{s = 0} \dot{\vec{q}}(t) + \cdots, t + s \at{\pderiv{T}{s}}{s = 0} + \cdots\right) \\
    &\quad\left(1 + s \deriv{}{t}\at{\pderiv{T}{s}}{s = 0} + \cdots\right) \d{t} \\
    &= \int_{t_A}^{t_B} \left[L(\vec{q}(t), \dot{\vec{q}}(t), t) - s \deriv{}{t}\at{\pderiv{T}{s}}{s = 0} \dot{q}_i \pderiv{L}{q_i} + s \at{\pderiv{T}{s}}{s=0} \pderiv{L}{t} \right. \\
      &\left.+ sL \deriv{}{t}\at{\pderiv{T}{s}}{s = 0} + O(s^2)\right] \d{t}
  \end{align*}
  (\textbf{TODO}: Finish)\par
  So $I^{*} = I$ thus $[\cdots] = 0\ \forall t_A, t_B$ and so:
  \[
    \at{\pderiv{T}{S}}{s = 0}\left(L - \dot{q}_i \pderiv{L}{\dot{q}_i}\right)
  \]
  is independent of $t$ and so we are done.
\end{proof}
\begin{example}
  Assume that $L$ has no explicit $t$ dependence, i.e $L(\vec{q}, \dot{\vec{q}})$.
  We set $T(t, s) = t + s$ and so $t^{*} = t + s$, note that this satisfies $T(t, 0) = t$.
  So then $\vec{q}^{*}(t^{*}) = \vec{q}(t) = \vec{q}(t^{*} - s)$ and thus:
  \[
    \deriv{\vec{q}^{*}}{{t^{*}}} = \dot{\vec{q}}(t^{*} - s)
  \]
  We then calculate:
  \begin{align*}
    \int_{t^{*}_A}^{t^{*}_B} L(\vec{q}^{*}(t^{*}), \deriv{\vec{q}^{*}}{{t^{*}}}(t^{*})) \d{t^{*}} &= \int_{t_A + s}^{t_B + s} L(\vec{q}(t^{*} - s), \dot{\vec{q}}(t^{*} - s)) \d{t^{*}} \\
                                                                                                  &= \int_{t_A}^{t_B} L(\vec{q}(t), \dot{\vec{q}}(t)) \d{x}
  \end{align*}
  by symmetry of the action.

  This has a conserved quantity:
  \[
    \underbrace{\at{\pderiv{T}{s}}{s = 0}}_{= 1}\left(L - \dot{q}_i \pderiv{L}{\dot{q}_i}\right) = - E
  \]
  Thus the invariance of the laws of physics under time transformations implies the conservation of energy.
\end{example}
\end{document}
